\chapter*{Preface}

\qquad Oftentimes, equations are given to students on various equation sheets on physics tests. Although knowing the equations is important, it is even more important to know the derivations of the equations as it provides a deeper understanding of the physical concepts behind the equations and builds physical intuition required to solve problems. Using the mathematical flexibility gained through proving formulas, students will better be able to apply their knowledge to solve more complex physics problems.
\par These notes attempt to derive as many formulae and concepts as possible that pertain to high school physics, from AP courses to Physics Olympiads. Equations featured in AP Physics 1 and 2 are derived without calculus (as much as possible), while calculus is employed when necessary to derive more difficult expressions. Vectors are represented using boldface ($\textbf{v}$) in this text, and, a star ($\bigstar$) is used to mark the conclusion of each concept/derivation.
\par Before we get into it, I also want to share my experience when I learned some of the concepts in these notes for the first time. I, like most other people on this planet, can't immediately conceptualize any topic that's thrown at me. I always takes multiple revisions over weeks or even months to actually get something into my brain. Take, for example, simple harmonic motion. When I first learned the topic, I could derive the period of a pendulum and solve for its period given a few parameters. When I reviewed the topic the second time, I became comfortable solving weirder systems, like a ball bobbing water or water sloshing in a U-tube. By the third revision (about an entire year after the first!) I finally understood damping and driven oscillations. All this is to say that learning is anything but linear. If you're stuck, take a break, come back, and oftentimes the thing you were stuck on becomes trivial. If you're still stuck, reach out to a teacher. If you don't have that luxury, I'm happy to help via email (physicshelprn63@gmail.com) or 1-on-1 tutor session. Just whatever you do, try to keep an open mind, even when you're frustrated at physics for being hard, at the problem you can't solve for being annoying, and at yourself for being stuck (trust me, anybody who learned a bit of physics as been through this).
\par Another thing to know about me: I'm not a math super genius. I'm passionate about physics, not math, which meant that while reading through harder concepts in some textbooks I found myself spending more time working through math than understanding physics! Thus, I've tried my best to show intermediate steps in mathematical derivations and transformations to make things a bit more readable. Still, you should derive everything by yourself after reading through, since you always learn more by doing :)
\par The sections required for various tests are marked with the following colors:
\begin{itemize}
\color{OrangeRed}
\item{$\blacktriangleright$} AP Physics 1
\color{RoyalBlue}
\item{$\blacktriangleright$} AP Physics 2
\color{PineGreen}
\item{$\blacktriangleright$} AP Physics C
\color{Goldenrod}
\item{$\blacktriangleright$} F=ma Exam
\color{Orchid}
\item{$\blacktriangleright$} USAPhO Exam
\end{itemize}

\mbox{}\linebreak
